% Chapter 5: Design

\section{Introduction}

This chapter presents the detailed design of the Lesaka Multi-Agent Orchestration System, including multi-agent architecture, graph orchestration design, agent design, state machine design, and self-healing mechanisms. The design phase translates the requirements identified in the analysis phase into concrete technical specifications.

\section{Multi-Agent Architecture}

\subsection{System Architecture Overview}

The multi-agent system follows a hierarchical architecture with three layers:

\begin{enumerate}
    \item \textbf{Orchestration Layer:} Graph orchestrator manages routing decisions
    \item \textbf{Agent Layer:} Specialized agents perform domain-specific analysis
    \item \textbf{Data Layer:} INFS database provides validated telemetry data
\end{enumerate}

\subsection{Component Architecture}

\begin{table}[H]
\centering
\begin{tabular}{|l|l|l|}
\hline
\textbf{Component} & \textbf{Responsibility} & \textbf{Interface} \\
\hline
Graph Orchestrator | Routing decisions & route\_telemetry() \\
Agent Molemo | Biomedical analysis | process() \\
Agent Loapi | Environmental analysis | process() |
Agent Thekiso | Market assessment | process() \\
Supervisor Agent | Error recovery | process() \\
\hline
\end{tabular}
\caption{Component Architecture}
\end{table}

\section{Graph Orchestration Design}

\subsection{Graph Structure}

The orchestration graph is modeled as a Directed Cyclic Graph (DCG) with the following structure:

\begin{itemize}
    \item \textbf{Nodes:} Agents (Molemo, Loapi, Thekiso, Supervisor)
    \item \textbf{Edges:} Routing paths based on temperature thresholds
    \item \textbf{Weights:} Priority levels for routing decisions
\end{itemize}

\subsection{Routing Algorithm}

The routing algorithm uses a state machine approach:

\begin{algorithm}
\caption{Telemetry Routing Algorithm}
\begin{algorithmic}[1]
\Require{telemetry\_data, cattle\_id}
\Ensure{agent\_response}
\State Fetch temperature from telemetry\_data
\State Validate context for cattle\_id
\If{context is invalid}
    \State Return Supervisor response
\EndIf
\If{temperature > 39.5}
    \State Route to Molemo (fever detection)
\ElsIf{temperature < 36.0}
    \State Route to Loapi (cold stress)
\Else
    \State Route to Thekiso (market assessment)
\EndIf
\State Return agent response
\end{algorithmic}
\end{algorithm}

\subsection{Algorithm Complexity}

\begin{itemize}
    \item \textbf{Time Complexity:} O(1) - Constant time for threshold comparison
    \item \textbf{Space Complexity:} O(1) - Constant space for routing decision
    \item \textbf{Determinism:} Deterministic routing based on temperature
\end{itemize}

\section{Agent Design}

\subsection{Agent Molemo Design}

\textbf{Specialization:} Biomedical analysis for fever detection

\textbf{Input:} Temperature, heart rate

\textbf{Output:} Severity level, recommendation, action required

\textbf{Thresholds:}
- Critical: $\geq$ 41.0°C
- High: $\geq$ 39.5°C
- Normal: < 39.5°C

\subsection{Agent Loapi Design}

\textbf{Specialization:} Environmental analysis for cold stress

\textbf{Input:} Temperature, humidity

\textbf{Output:} Severity level, recommendation, action required

\textbf{Thresholds:}
- Severe: $\leq$ 34.0°C
- Moderate: $\leq$ 36.0°C
- Normal: > 36.0°C

\subsection{Agent Thekiso Design}

\textbf{Specialization:} Market assessment for liquidation timing

\textbf{Input:} Temperature, heart rate

\textbf{Output:} Market ready status, grade, estimated value

\textbf{Grading:}
- Grade A: Optimal health
- Grade B: Market-ready with considerations
- Grade C: Requires health assessment
- Grade D: Not suitable for market

\subsection{Supervisor Agent Design}

\textbf{Specialization:} Error recovery and system stability

\textbf{Input:} Cattle ID, error type, error message

\textbf{Output:} Fallback action, recovery message

\textbf{Fallback Strategies:}
- No data found: Request data refresh
- Context anomaly: Re-validate regional identifier
- Processing error: Retry in safe mode
- Unknown error: Escalate to manual review

\section{State Machine Design}

\subsection{State Definition}

The state machine has the following states:

\begin{table}[H]
\centering
\begin{tabular}{|l|l|l|}
\hline
\textbf{State} & \textbf{Condition} & \textbf{Action} \\
\hline
Normal | 36.0°C $\leq$ temp $\leq$ 39.5°C | Route to Thekiso \\
Fever | temp > 39.5°C | Route to Molemo \\
Cold Stress | temp < 36.0°C | Route to Loapi \\
Error | Any error condition | Route to Supervisor \\
\hline
\end{tabular}
\caption{State Machine Definition}
\end{table}

\subsection{State Transitions}

State transitions are deterministic based on temperature thresholds:

\begin{itemize}
    \item Normal $\rightarrow$ Fever: Temperature exceeds 39.5°C
    \item Normal $\rightarrow$ Cold Stress: Temperature drops below 36.0°C
    \item Fever $\rightarrow$ Normal: Temperature returns to normal range
    \item Cold Stress $\rightarrow$ Normal: Temperature returns to normal range
    \item Any $\rightarrow$ Error: Processing error detected
    \item Error $\rightarrow$ Normal: Error resolved
\end{itemize}

\section{Self-Healing Design}

\subsection{Error Detection}

The system detects errors through:

\begin{itemize}
    \item Exception handling in routing logic
    \item Context validation before routing
    \item Data availability checks
    \item Agent response validation
\end{itemize}

\subsection{Fallback Strategies}

The system implements the following fallback strategies:

\begin{table}[H]
\centering
\begin{tabular}{|l|l|l|}
\hline
\textbf{Error Type} & \textbf{Fallback Action} & \textbf{Recovery Message} \\
\hline
No Data Found | Request Data Refresh | Requesting data refresh \\
Context Anomaly | Validate Context | Re-validating regional identifier \\
Processing Error | Retry Safe Mode | Retrying in safe mode \\
Unknown Error | Escalate | Escalating to manual review \\
\hline
\end{tabular}
\caption{Fallback Strategies}
\end{table}

\section{Integration Design}

\subsection{API Contract with INFS Project}

The Lesaka Multi-Agent Orchestration System consumes data from the INFS Data Governance Framework through a well-defined API contract:

\begin{itemize}
    \item \textbf{Data Access:} Query INFS database for validated telemetry
    \item \textbf{Permission Respect:} Respect INFS RBAC restrictions
    \item \textbf{Error Handling:} Standardized error responses
    \item \textbf{Data Format:} JSON format for data exchange
\end{itemize}

\subsection{Integration Points}

\begin{itemize}
    \item \textbf{Database Layer:} Query INFS database for telemetry
    \item \textbf{Validation Layer:} Use INFS-validated data
    \item \textbf{Security Layer:} Respect INFS RBAC permissions
\end{itemize}

\section{Deployment Design}

\subsection{Development Environment}

\begin{itemize}
    \item Local development with virtual environment
    \item Python 3.10+ runtime
    \item NetworkX for graph operations
    \item SQLite database (provided by INFS project)
\end{itemize}

\subsection{Production Considerations}

Future production deployment would include:

\begin{itemize}
    \item Enhanced error logging
    \item Performance monitoring
    \item Load balancing for high-volume routing
    \item Distributed agent deployment
\end{itemize}

\section{Conclusion}

The design phase has established a comprehensive technical blueprint for the Lesaka Multi-Agent Orchestration System. The design addresses the requirements identified in the analysis phase while incorporating best practices for multi-agent systems, graph algorithms, and self-healing mechanisms. The design is ready for implementation in the next phase.
